\section{Mengapa Teknik Kompilasi Penting}
\label{sec:ch01-mengapa}

Teknik Kompilasi adalah mata kuliah yang menjelaskan bagaimana sebuah
program sumber berbahasa tingkat tinggi diterjemahkan menjadi bentuk
yang dapat dijalankan oleh mesin. Pemahaman ini penting bukan hanya
untuk mahasiswa yang ingin membangun kompilator, tetapi juga untuk
memahami alat bantu pengembangan, bahasa domain khusus, analisis kode,
dan sistem kecerdasan buatan yang memproses bahasa.

Pekerjaan seorang insinyur perangkat lunak sangat sering bersinggungan
dengan teknologi yang lahir dari teknik kompilasi: \eng{parser} pada
editor, \eng{linter} dan \eng{static analyzer}, mesin \eng{template},
\eng{query optimizer} basis data, hingga penerjemah bahasa alami.
Kemampuan membaca spesifikasi bahasa, memodelkan tata bahasa, dan
menganalisis kesalahan secara sistematis adalah kompetensi yang dapat
dipakai di banyak bidang.

Buku ini mengikuti pendekatan \eng{Outcome-Based Education} (OBE).
Setiap bab dimulai dengan outcome yang harus dicapai, dilanjutkan
dengan materi, contoh, praktikum, latihan, asesmen, dan daftar cek
kompetensi. Dengan cara ini, mahasiswa tidak hanya membaca teori,
tetapi juga dapat menunjukkan bukti pencapaian kompetensi
\citep{thain2024}.

\begin{infobox}
Mata kuliah ini menekankan \textbf{front-end} kompilator: analisis
leksikal, analisis sintaksis, analisis semantik, serta pembentukan dan
validasi AST. Optimasi dan pembangkitan kode dibahas sebagai konteks
pada Bab III dan Bab XVII, bukan sebagai bab tersendiri.
\end{infobox}
